\section{The Fix: Structural, Differentiable Orthogonalization (F-geo-3)}
\label{sec:the-fix}

% Full grounding: DELTANET_RD_EXACTNESS_DESIGN.md \S14 (esp. \S14.0, \S14.1,
% \S14.6, \S14.10).

\subsection{Motivation (the two facts, not engineering for its own sake)}
\begin{itemize}
  \item (1) The exact solution exists and composes perfectly (i-strong,
    \S\ref{sec:mechanism}).
  \item (2) SGD does not find it under any soft push tried
    (\S\ref{sec:soft-fixes}).
\end{itemize}

\subsection{Design}
\begin{itemize}
  \item Cubic Newton-Schulz / Bj\"orck-Bowie iteration at the
    $k_{\mathrm{eff}}$ site: $X_{t+1} = 1.5 X_t - 0.5 (X_t X_t^\top) X_t$,
    converges quadratically to the nearest orthonormal-row matrix given
    operator norm $\leq 1$ at $t{=}0$ (guaranteed by $\sqrt{K}$
    pre-scaling given the existing unit-row-norm invariant).
  \item REJECTED: Gram-Schmidt (order-dependent --- would reintroduce
    exactly mechanism (g)'s write-history asymmetry one level up); QR
    (order-dependent for the same underlying reason, AND backward-unstable
    near collinear input --- the exact regime an undertrained embedding
    table produces).
  \item Chosen properties: order-equivariant (permuting input rows
    permutes iterates identically --- no row privileged), differentiable
    through ordinary matmuls (no LinAlg-specific backward formula).
\end{itemize}

\subsection{The gating discipline (concrete instance 2, per \S\ref{sec:setup})}
\begin{itemize}
  \item Two-round independent attack found a load-bearing failure mode
    the design missed: cross-episode key-identity DRIFT ---
    orthogonalizing correctly WITHIN an episode does not guarantee the
    SAME key across episodes of the same entity.
  \item Became a pre-registered Wave $-1$ gating diagnostic with its own
    committed simulator (\texttt{geo3\_simulator.py}, importable,
    CPU-smoke-tested) and launch-read rule.
  \item Wave 1 GPU spend authorized ONLY after the simulator predicted the
    $K{=}16$ bar would clear: predicted $K16\ h4{=}1.00$, $K32\ h4{=}0.77$
    from measured drift $0.94/0.90$.
  \item Framing rule for the actual prose: pair every process claim with
    the SPECIFIC finding it caught (as above) --- never narrate ``we used
    a rigorous review process'' in the abstract.
\end{itemize}

\subsection{The substitute admission stack (now footnote-tier, but shown,
not asserted)}
\begin{itemize}
  \item F-geo-3 makes the standard key-Gram/alignment premise-validity
    instruments TAUTOLOGICAL by construction (orthogonality $\equiv$
    true, alignment $\equiv 1.0$).
  \item Substitute, pre-registered BEFORE any F-geo-3 data existed: value-
    side salvage tier + Newton-Schulz convergence WITHOUT eigh fallback +
    finite loss + a task-performance floor ($\geq$50\% loss improvement
    AND $h{=}1$ ID $\mathrm{rec@0.9} \geq 0.5$).
  \item \textbf{Discharge status (resolved):} with the $K{=}32$
    \texttt{n\_iter=20} escalation cells passing EVERY substitute
    criterion including zero fallback steps (and $K{=}16$ 3/3 from the
    start), this caveat drops from headline asterisk to footnote --- but
    the footnote still names both admission stacks and both realized pass
    rates wherever geo3 and non-geo3 seed counts appear side by side, per
    the pre-registration's own requirement. The fallback-irrelevance
    finding (\S\ref{sec:results}) shows the one criterion that ever
    disqualified a seed was conservative, not explanatory.
\end{itemize}

\subsection{The direct follow-on to the outcome-F residual --- designed,
attacked, built, run to completion across four waves, and now CLOSED
(Fig.~10, Fig.~11)}
\begin{itemize}
  \item \textbf{This is no longer future work.} The outcome-F residual
    (\S\ref{sec:results}: cross-episode key drift $0.9037$ in the
    registered HIGH band is the $K{=}32$ bar-miss cause) motivated a
    trainable, per-entity anchor table blended into the key before geo3's
    Newton-Schulz pass. The full arc, in the order it happened:
  \item \textbf{(1) Behavioral gain, launched twice, then independently
    replicated.} Initial + confirmatory waves produced a real, seed-stable
    $K{=}32$ $h{=}4$ gain (freshly-measured baseline $0.4105 \to
    0.556$--$0.665$). A mechanism-tier wave (6 genuinely fresh seeds, 2
    architecture variants) reconfirmed the gain at $0.6141$--$0.7141$ ---
    9 total seed-runs across 3 distinct seed sets license
    ``independently replicated,'' not merely ``reproduced.''
  \item \textbf{(2) The hypothesized mechanism was tested and REJECTED
    (Outcome C).} The confirmatory wave's own per-entity engagement
    readout showed $<$14\% engagement, not Outcome A. A same-day rescue
    attempt (a $\lambda$-scaled engagement-threshold re-derivation) was
    independently attacked and rejected --- reported as attempted and
    rejected, not deleted. The mechanism-tier wave reconfirmed Outcome C at
    a strictly stronger evidentiary bar (unit-anchor-norm measured, not
    assumed, at $1.06$--$1.16$; a formal dip-test rules out a masked
    bimodal split; an effect-size floor rules out power-manufactured
    engagement) and registered a falsifiable alternative account BEFORE
    the confirming data existed: ``cross-episode stability comes from the
    mere presence of an episode-constant blend term, not learned
    alignment.''
  \item \textbf{(3) That prediction was tested and HELD ---
    CONFIRMED-BY-ABLATION.} Candidate (e): two frozen, never-trained
    anchor tables (pure-random-init and frame-potential-init) at the same
    fixed blend weight ($\lambda{=}0.58$) both match or slightly exceed
    the learned table's own $K{=}32$ mean ($0.7274$/$0.7413$ vs.\ $(d)$'s
    $0.6669$; no frozen seed falls below $(d)$'s own minimum). The $r_e$
    negative control passes at the strongest null in the design's history
    (median $r_e$ negative, $-0.13$ to $-0.24$, for the pure-random arm,
    which cannot mechanically align to anything). Constancy alone
    suffices; ``learned anchoring'' is superseded by a frozen random
    key-bias at matched blend weight, no training loop required.
  \item \textbf{(4) The fix's operating regime was then measured and
    bounded, not left open-ended.} A capacity-curve extension to $K{=}48$
    ($K/d{=}0.75$): candidate (d) fully admissible (3/3 seeds), both named
    explanatory channels (drift-space closure, $\sim$81\% of the gap;
    value-Gram relief, $0.626\times$ the reference) measurably active, bar
    still missed 0/3 seeds (mean $0.0215$ vs.\ $0.024434$). The
    three-point curve ($\sim$1.00 $\to$ $\sim$0.65 $\to$ $\sim$0.02 at
    $K/d{=}0.25/0.5/0.75$) is a cliff, not a smooth decline, while the
    theoretical $\lambda{=}1$ ceiling declines only gently over the same
    range ($0.97 \to 0.94 \to 0.90$).
  \item \textbf{Program status: COMPLETE.} Five waves plus one rejected
    rescore attempt, $\approx$55.83/80 GPU-h against its own budget cap, no
    further waves scheduled. The one genuinely open follow-on this program
    itself named (value-crowding at $K/d{=}0.75$) is a NEW, unregistered
    direction requiring its own fresh design/attack/build/audit cycle ---
    it is not a continuation of this design.
  \item Frame sentence for this section: \emph{a rejected hypothesis and a
    confirmed-by-ablation alternative are both reported, not just the
    flattering one --- and the alternative's own operating regime is now
    bounded, closing what was, at an earlier stage of this project, this
    paper's single most open thread.}
\end{itemize}
